\section{Target Localisation}\label{sec:localisation}

Once the onboard detector identifies a candidate target in an image frame, the system must convert that pixel-level observation into a geographic coordinate suitable for navigation. A single detection is inherently noisy---subject to GPS receiver latency, altitude uncertainty, and finite pixel resolution---so the system fuses multiple observations over time using spatial clustering, inverse-variance weighting, and Kalman filtering.

% ─────────────────────────────────────────────────────────────────────
\subsection{Pixel-to-GPS Conversion}\label{sec:pixel-gps}

Each detection returns the bounding-box centre $(u, v)$ in pixel coordinates. The conversion to a ground-plane offset uses the pinhole camera model with the following parameters: sensor width $w_s = \SI{5.02}{\milli\metre}$ (IMX296), calibrated focal length $f = \SI{5.46}{\milli\metre}$, and image dimensions $W \times H = 1456 \times 1088$~pixels. The ground sampling distance (GSD) at altitude $h$ is

\begin{equation}\label{eq:gsd}
  \mathrm{GSD} = \frac{w_s \cdot h}{f \cdot W} \qquad [\text{m\,px}^{-1}]
\end{equation}

\noindent At the nominal search altitude of $h = \SI{35}{\metre}$, this yields $\mathrm{GSD} \approx 0.022$\,m\,px$^{-1}$, so the full frame covers roughly $\SI{32}{\metre} \times \SI{24}{\metre}$ on the ground.

The pixel offset from the image centre $(C_x, C_y) = (W/2,\; H/2)$ is projected to a forward--right body-frame displacement:

\begin{equation}\label{eq:body-offset}
  d_\mathrm{fwd} = -(v - C_y)\,\mathrm{GSD}, \qquad
  d_\mathrm{right} = (u - C_x)\,\mathrm{GSD}
\end{equation}

\noindent These are rotated into the North--East frame using the drone's heading $\psi$:

\begin{equation}\label{eq:ne-offset}
  \Delta N = d_\mathrm{fwd}\cos\psi - d_\mathrm{right}\sin\psi, \qquad
  \Delta E = d_\mathrm{fwd}\sin\psi + d_\mathrm{right}\cos\psi
\end{equation}

\noindent Finally, the offsets are converted to latitude and longitude increments on the WGS\,84 ellipsoid:

\begin{equation}\label{eq:gps-offset}
  \Delta\phi = \frac{\Delta N}{R_\oplus}\cdot\frac{180}{\pi}, \qquad
  \Delta\lambda = \frac{\Delta E}{R_\oplus \cos\phi}\cdot\frac{180}{\pi}
\end{equation}

\noindent where $R_\oplus = \SI{6378137}{\metre}$ is the WGS\,84 semi-major axis and $\phi$ is the current latitude. The estimated target position is then $(\phi + \Delta\phi,\; \lambda + \Delta\lambda)$.

When the target falls within a small threshold of the image centre (the ``centre-snap'' zone), the pixel projection is bypassed and the drone's own GPS position is used directly, since the target is approximately directly below the camera.

% ─────────────────────────────────────────────────────────────────────
\subsection{Multi-Detection Fusion}\label{sec:fusion}

A single GPS estimate is insufficient for reliable landing guidance. The system therefore accumulates observations across successive frames and applies two fusion mechanisms.

\paragraph{Spatial clustering.}
Each new observation is compared against existing detection clusters using the Haversine distance. If the nearest cluster centroid is within a configurable threshold (default \SI{30}{\metre}), the observation is added to that cluster; otherwise, a new cluster is created. This allows the system to track multiple candidate targets simultaneously---for example, a real casualty and a piece of debris---without conflating their position estimates. Each cluster maintains its own rolling window of the most recent 50 observations and an all-time running average.

\paragraph{Inverse-variance weighting.}
Not all observations are equally reliable. Two factors modulate the weight $w_i$ assigned to each observation:

\begin{enumerate}
  \item \textbf{Altitude factor.} Position error scales linearly with altitude (via the GSD), so variance scales quadratically. The weight is therefore $w_\mathrm{alt} = (h_\mathrm{ref} / h)^2$, with $h_\mathrm{ref} = \SI{30}{\metre}$. At \SI{15}{\metre} this gives a $4\times$ weight increase; at \SI{10}{\metre}, $9\times$.
  \item \textbf{Centrality factor.} Detections near the image centre suffer less from lens distortion and heading uncertainty. The weight is linearly interpolated from 5 (at centre) to 1 (at corners):
  $w_\mathrm{cen} = 1 + 4\,(1 - d / d_\mathrm{max})$, where $d$ is the pixel distance from the image centre and $d_\mathrm{max}$ is the diagonal half-length.
\end{enumerate}

\noindent The combined weight is $w_i = w_\mathrm{cen} \cdot w_\mathrm{alt}$, and the fused cluster position is the weighted mean:

\begin{equation}\label{eq:weighted-avg}
  \hat{\phi} = \frac{\sum_i w_i\,\phi_i}{\sum_i w_i}, \qquad
  \hat{\lambda} = \frac{\sum_i w_i\,\lambda_i}{\sum_i w_i}
\end{equation}

% ─────────────────────────────────────────────────────────────────────
\subsection{Kalman Filtering}\label{sec:kalman}

While the weighted average converges given enough observations, it treats all recent measurements equally (within the 50-sample window) and has no formal model of uncertainty. A two-state Kalman filter~\cite{kalman1960new} is therefore applied in parallel, operating on the latitude--longitude state vector $\mathbf{x}_k = [\phi,\;\lambda]^\top$.

The target is assumed stationary (casualty on the ground), so the state transition is the identity: $\mathbf{x}_{k|k-1} = \mathbf{x}_{k-1}$. A small process noise $\mathbf{Q} = \sigma_q^2\,\mathbf{I}$ is applied, where $\sigma_q$ corresponds to \SI{0.1}{\metre} of positional uncertainty converted to degrees via the WGS\,84 semi-major axis ($\sigma_q \approx 0.1 / R_\oplus \cdot 180/\pi$). This accounts for minor model mismatch. The measurement noise covariance is set inversely proportional to the observation weight:

\begin{equation}\label{eq:kalman-R}
  \sigma_R = \frac{3\;\text{m}}{\max(w_i,\; 0.1)}, \qquad
  \mathbf{R}_k = \sigma_R^2\,\mathbf{I}
\end{equation}

\noindent so that high-weight observations (low altitude, near image centre) pull the state estimate more aggressively than high-altitude edge detections. The standard predict--update cycle follows:

\begin{align}
  \mathbf{P}_{k|k-1} &= \mathbf{P}_{k-1} + \mathbf{Q} \label{eq:kalman-predict}\\
  \mathbf{K}_k &= \mathbf{P}_{k|k-1}\,(\mathbf{P}_{k|k-1} + \mathbf{R}_k)^{-1} \label{eq:kalman-gain}\\
  \mathbf{x}_k &= \mathbf{x}_{k|k-1} + \mathbf{K}_k\,(\mathbf{z}_k - \mathbf{x}_{k|k-1}) \label{eq:kalman-update}\\
  \mathbf{P}_k &= (\mathbf{I} - \mathbf{K}_k)\,\mathbf{P}_{k|k-1} \label{eq:kalman-cov}
\end{align}

\noindent where $\mathbf{z}_k = [\phi_\mathrm{est},\;\lambda_\mathrm{est}]^\top$ is the pixel-derived GPS measurement. The filter converges within approximately 5--10 observations, after which the covariance $\mathbf{P}$ stabilises and additional measurements produce diminishing corrections. In simulation, the Kalman estimate consistently outperforms the simple weighted average once more than 10 observations are accumulated.

% ─────────────────────────────────────────────────────────────────────
\subsection{GPS Error Sources}\label{sec:gps-error}

Several factors limit the achievable localisation accuracy:

\begin{itemize}
  \item \textbf{Receiver latency.} Consumer-grade GPS receivers exhibit \SIrange{100}{200}{\milli\second} of position latency~\cite{misra2006global}. At a search speed of \SI{5}{\metre\per\second}, this introduces up to \SI{1}{\metre} of along-track bias in every observation.
  \item \textbf{GPS multipath and drift.} Even with a clear sky view, single-frequency L1 receivers have a circular error probable (CEP) of \SIrange{2}{3}{\metre}~\cite{kaplan2017understanding}. This noise is largely isotropic and averages out with sufficient observations, but individual estimates can deviate by \SI{5}{\metre} or more.
  \item \textbf{Altitude uncertainty.} Barometric altitude is used for the GSD calculation. Barometric drift of $\pm\SI{1}{\metre}$ at \SI{35}{\metre} induces a ${\sim}3\%$ GSD error, shifting the ground projection by up to \SI{0.5}{\metre} for edge-of-frame detections.
  \item \textbf{Heading uncertainty.} The rotation from body frame to North--East (Equation~\ref{eq:ne-offset}) depends on the magnetometer heading. A $\pm\SI{5}{\degree}$ compass error at \SI{35}{\metre} altitude displaces the estimate by up to \SI{1.4}{\metre} for targets at the image edge.
  \item \textbf{FOV calibration error.} The focal length was calibrated to $f = \SI{5.46}{\milli\metre}$ by measuring the visible width at a known distance (Section~\ref{sec:pixel-gps}). Residual calibration error of $\pm5\%$ scales all ground offsets proportionally.
\end{itemize}

\noindent GPS latency is the dominant systematic error during forward flight, as it consistently biases estimates in the direction of travel. This is partially mitigated by the multi-pass nature of the lawnmower pattern: alternating scan directions cause the latency bias to cancel in the weighted average.

% ─────────────────────────────────────────────────────────────────────
\subsection{Accuracy Analysis}\label{sec:accuracy}

The localisation pipeline was validated by replaying recorded DJI flight video ($1456 \times 1088$~pixels at 30\,fps) with synchronised SRT telemetry containing per-frame GPS and altitude data. A known dummy target was placed at a surveyed ground-truth position, and each detection frame produced an independent GPS estimate. The results are summarised in Table~\ref{tab:localisation-accuracy}.

\begin{table}[ht]
  \centering
  \caption{Target localisation accuracy from DJI video replay (\SI{54.4}{\degree} HFOV, \SIrange{15}{50}{\metre} altitude range).}
  \label{tab:localisation-accuracy}
  \begin{tabular}{lc}
    \toprule
    \textbf{Metric} & \textbf{Value} \\
    \midrule
    CEP50 (50\% of estimates within) & \SI{2.3}{\metre} \\
    CEP95 (95\% of estimates within) & \SI{8.5}{\metre} \\
    Maximum error (single observation) & \SI{16.5}{\metre} \\
    Weighted-average error (all obs.) & $<$\,\SI{2}{\metre} \\
    Kalman-filtered error (converged) & $<$\,\SI{1.5}{\metre} \\
    \bottomrule
  \end{tabular}
\end{table}

\begin{figure}[ht]
  \centering
  % \includegraphics[width=0.7\columnwidth]{figures/gps_scatter.pdf}
  \textit{[Placeholder: GPS scatter plot showing individual estimates (grey dots), weighted average (blue cross), Kalman estimate (red cross), and ground truth (green circle). Axes in metres relative to ground truth.]}
  \caption{Scatter plot of target position estimates from video replay. Individual observations (grey) are dispersed along the flight direction due to GPS latency; the fused estimates (blue, red) converge to within \SI{2}{\metre} of ground truth.}
  \label{fig:gps-scatter}
\end{figure}

The scatter plot (Figure~\ref{fig:gps-scatter}) shows the characteristic elongation along the flight direction caused by GPS receiver latency. The Kalman filter and weighted average both converge to within approximately \SI{2}{\metre} of the true position after 10--15 observations. This accuracy is sufficient for the approach phase, where the drone descends and switches to visual servoing (pixel-based centring) for the final \SI{15}{\metre} of altitude.

Two strategies could further reduce localisation error in future work: (i)~compensating for GPS latency by subtracting $v_{\mathrm{ground}} \cdot \Delta t$ along the heading vector, where $\Delta t \approx \SI{150}{\milli\second}$ is the estimated receiver delay; and (ii)~fusing RTK-corrected GPS, which would reduce the base position noise from approximately \SI{2.5}{\metre} CEP to below \SI{0.02}{\metre}~\cite{takasu2009development}.
